% Geometry of the two-variable LP: max 2 z1 + 3 z2 s.t. z1 + z2 <= 4, z1 + 3 z2 <= 6, z >= 0.
\begin{tikzpicture}[font=\footnotesize,x=1.25cm,y=1.25cm]
  \fill[sbregionB,draw=none] (0,0) -- (4,0) -- (3,1) -- (0,2) -- cycle;
  \draw[->,black!60] (-0.3,0) -- (6.6,0) node[below] {$z_1$};
  \draw[->,black!60] (0,-0.3) -- (0,4.4) node[left] {$z_2$};
  \foreach \t in {1,...,6} \draw[black!60] (\t,-0.06) -- (\t,0.06) node[below,yshift=-2pt] {\t};
  \foreach \t in {1,...,4} \draw[black!60] (-0.06,\t) -- (0.06,\t) node[left,xshift=-2pt] {\t};
  % constraint lines
  \draw[sbBlue,thick] (0,4) -- (4.3,-0.3) node[sbannot,text=sbBlue,right,xshift=-2pt,yshift=4pt] {$z_1+z_2=4$};
  \draw[sbBlue,thick] (-0.3,2.1) -- (6.3,-0.1) node[sbannot,text=sbBlue,above right,xshift=-8pt] {$z_1+3z_2=6$};
  % level lines of the objective 2 z1 + 3 z2 = c
  \draw[sbOrange,densely dotted] (1.5,0) -- (0,1) node[sbannot,text=sbOrange,left] {$3$};
  \draw[sbOrange,densely dotted] (3,0) -- (0,2) node[sbannot,text=sbOrange,left,yshift=6pt] {$6$};
  \draw[sbOrange,densely dotted] (4.5,0) -- (0,3) node[sbannot,text=sbOrange,left] {$9$};
  \draw[sbOrange,densely dotted] (6,0) -- (0.9,3.4) node[sbannot,text=sbOrange,above] {$c=12$};
  % vertices and their objective values
  \fill[black] (0,0) circle (1.5pt) node[below left,sbannot,text=black] {$0$};
  \fill[black] (4,0) circle (1.5pt) node[below,sbannot,text=black,yshift=-3pt] {$8$};
  \fill[black] (0,2) circle (1.5pt) node[above left,sbannot,text=black] {$6$};
  \node[sbgoal] at (3,1) {};
  \node[sbannot,text=black,below right,yshift=-2pt] at (3,1) {$\vect{z}^*=(3,1)$, value $9$};
  % objective gradient and active normals at the optimum
  \draw[sbvec,draw=sbOrange] (3,1) -- (3.5,1.75) node[sbannot,text=sbOrange,above] {$\vect{c}=(2,3)$};
  \draw[sbvec,draw=sbBlue] (3,1) -- (3.55,1.55) node[sbannot,text=sbBlue,right] {$(1,1)$};
  \draw[sbvec,draw=sbBlue] (3,1) -- (3.25,1.75) node[sbannot,text=sbBlue,left,xshift=1pt] {$(1,3)$};
  \node[sbannot,text=black] at (1.2,0.55) {$P$};
\end{tikzpicture}
